% Ad Atticum 13.23 (ad-atticum-13-23)
% Source: https://raw.githubusercontent.com/PerseusDL/canonical-latinLit/master/data/phi0474/phi057/phi0474.phi057.perseus-lat1.xml
% Fetched by scripts/fetch_latin.py

% Dateline (Perseus): Scr. in Tusculano vi Id. Quint. a. 709 (45).

\ciceroLetterOpener{CICERO ATTICO salutem}

\ciceroSection{1} ante meridianis tuis litteris heri statim rescripsi; nunc respondeo vespertinis. Brutus mallem me arcesseret. et aequius erat, cum illi iter instaret et subitum et longum, et me hercule nunc, cum ita simus adfecti ut non possimus plane simul vivere (intellegis enim profecto in quo maxime posita sit συμβίωσισ ), facile patiebar nos potius Romae una esse quam in Tusculano.

\ciceroSection{2} libri ad Varronem non morabantur, sunt enim †deffecti†, ut vidisti; tantum librariorum menda tolluntur. de quibus libris scis me dubitasse, sed tu videris. item quos Bruto mittimus in manibus habent librarii.

\ciceroSection{3} mea mandata, scribis, explica. quamquam ista retentione omnis ait uti Trebatius; quid tu istos putas? nosti domum. qua re confice εὐγαγώγωσ . incredibile est quam ego ista non curem. omni tibi adseveratione adfirmo, quod mihi credas velim, mihi maiori offensioni esse quam delectationi possessiunculas meas. magis enim doleo me non habere quoi tradam quam habere qui utar laetor . atque illud Trebatius se tibi dixisse narrabat; tu autem veritus es fortasse ne ego invitus audirem. fuit id quidem humanitatis, sed, mihi crede, iam ista non curo. qua re da te in sermonem et perseca et confice et ita cum Polla loquere ut te cum illo Scaeva loqui putes nec existimes eos qui non debita consectari soleant quod debeatur remissuros. de die tantum videto et id ipsum bono modo.
